\documentclass[letterpaper]{article}

\usepackage[utf8]{inputenc}
\usepackage[T1]{fontenc}
\usepackage{amsmath,amssymb}
\usepackage{booktabs}
\usepackage{multirow}
\usepackage{graphicx}
\usepackage{xcolor}
\usepackage[margin=1in]{geometry}
\usepackage{hyperref}
\usepackage{natbib}
\usepackage{caption}
\usepackage{subcaption}

\title{What Does Fine-Tuning Actually Learn? \\Representational Analysis of Action Structure in Vision-Language-Action Models}

\author{
  Anonymous Authors
}

\date{}

\begin{document}
\maketitle

%============================================================================
% ABSTRACT
%============================================================================
\begin{abstract}
Vision-Language-Action models (VLAs) achieve strong robotic manipulation performance by fine-tuning large language models on vision-conditioned action prediction.
A widespread assumption is that this fine-tuning process injects new action-relevant knowledge into the model's representations.
We challenge this assumption through a systematic representational and functional analysis of an OpenVLA-class architecture, decomposing contributions of the frozen vision encoder (SigLIP), projector, and fine-tuned LLM backbone.
Our analysis reveals three findings on LIBERO-Goal, with partial replication on LIBERO-Spatial.
\textbf{First}, VL-pretrained vision encoders provide representations with strong geometric alignment to action structure (SigLIP RSA $\rho{=}0.615$), substantially exceeding the full VLA's LLM layers ($\rho{\approx}0.50$) and a self-supervised baseline (DINOv2, $\rho{=}0.388$); a CLIP-ViT control ($\rho{=}0.46$) confirms that SigLIP's specific pretraining recipe confers additional alignment beyond generic VL pretraining.
\textbf{Second}, the LLM degrades distributed geometric alignment (RSA: $0.615 {\to} 0.501$) while preserving global structure (CKA increases monotonically across layers), and training-dependently concentrates action information into a readout token (trained $R^2{=}0.731$ vs.\ untrained $0.686$, $\Delta{=}0.045$, $p{<}0.0001$).
\textbf{Third}, distributed representation structure is training-invariant (trained vs.\ untrained image-mean RSA: $0.501$ vs.\ $0.485$, bootstrap CI includes~0, $p{=}0.20$) and instruction-invariant ($\Delta{<}0.017$ across four instruction conditions), while readout concentration is training-dependent.
We conclude that VLA fine-tuning learns \emph{information routing}---how to channel pre-existing action-relevant structure to the output position---rather than creating new action knowledge.
\end{abstract}

%============================================================================
% 1. INTRODUCTION
%============================================================================
\section{Introduction}

Vision-Language-Action models (VLAs) have emerged as a promising paradigm for robotic manipulation, leveraging the representational capacity of large vision-language models to predict actions from visual observations and language instructions \citep{brohan2023rt2,kim2024openvla,karamcheti2024prismatic}.
The standard recipe---freezing a pretrained vision encoder, projecting visual tokens into a language model, and fine-tuning the LLM on action-labeled robot data---has yielded impressive generalization.
Yet a fundamental question remains open: \emph{where does the action-relevant representational structure come from?}

The implicit assumption in the VLA literature is that task-specific fine-tuning injects new action knowledge into the LLM's representations.
Recent behavioral studies suggest otherwise: Grant et al.\ (2026) find that the visual pathway dominates action selection, while Hou (2026) and Fei et al.\ (2025) independently report that VLAs largely ignore language instructions.
These behavioral observations call for a \emph{representational-level} explanation.

We provide the first systematic representational and functional decomposition of an OpenVLA-class VLA, separately quantifying the contributions of the frozen vision encoder (SigLIP), the linear projector, and the fine-tuned LLM backbone.
Our analysis combines Representational Similarity Analysis (RSA), Centered Kernel Alignment (CKA), and functional action probes across a $2{\times}2$ design (trained/untrained $\times$ with/without instruction) plus encoder controls (CLIP-ViT, DINOv2).

We find that the frozen SigLIP encoder already provides representations with the strongest geometric alignment to action structure (RSA $\rho{=}0.615$), exceeding any LLM layer.
The LLM degrades this alignment at the distributed image-token level but learns to \emph{concentrate} action-relevant information into the readout token---a training-dependent routing operation (readout $\Delta R^2{=}0.045$, $p{<}0.0001$) that leaves distributed structure unchanged ($p{=}0.20$).
Language instructions affect neither distributed nor concentrated representations.
We thus reframe VLA fine-tuning as \emph{routing without creating}: the LLM learns to channel pre-existing visual action structure to the output position, rather than building new action representations.

Our contributions are:
\begin{enumerate}
  \item A component-level decomposition showing VL-pretrained vision encoders dominate action-relevant geometric alignment, with SigLIP's specific recipe providing additional advantage beyond generic VL pretraining (CLIP-ViT control: RSA $0.46$).
  \item Evidence that the LLM selectively degrades local geometric structure while preserving global structure, with a training-dependent readout concentration mechanism ($\Delta{=}0.045$, $p{<}0.0001$).
  \item A $2{\times}2$ factorial analysis establishing that distributed representation structure is both training- and instruction-invariant, reframing fine-tuning as information routing.
\end{enumerate}

%============================================================================
% 2. RELATED WORK
%============================================================================
\section{Related Work}

\paragraph{Vision-Language-Action Models.}
VLAs build on vision-language models (VLMs) by fine-tuning the LLM component for action prediction \citep{brohan2023rt2,kim2024openvla}.
Prismatic VLMs \citep{karamcheti2024prismatic} provide the architectural foundation for OpenVLA, using SigLIP as the vision encoder with a linear projector into a Llama-class LLM.
Our work complements this line by analyzing what fine-tuning actually changes in the representational structure.

\paragraph{Behavioral Analysis of VLAs.}
Grant et al.\ (2026) demonstrate that the visual pathway dominates VLA action selection at the behavioral level.
Hou (2026; LangGap) and Fei et al.\ (2025; LIBERO-Plus) independently find that VLAs largely ignore language instructions behaviorally.
We provide a representational-level explanation: instruction content has negligible effect on internal representations ($\Delta \text{RSA} < 0.017$).

\paragraph{Mechanistic Interpretability of VLAs.}
Haon et al.\ (2025) apply activation steering to find sparse semantic directions in OpenVLA.
Swann et al.\ (2026) use sparse autoencoders for OpenVLA interpretability.
Lu et al.\ (2025) show that linear probes on OpenVLA hidden layers achieve $>90\%$ accuracy on semantic features.
Molinari et al.\ (2025) study emergent world representations.
Our geometric analysis complements these approaches by quantifying how action-relevant \emph{structure}---not individual features---is transformed across components.

\paragraph{Representational Analysis Methods.}
RSA \citep{kriegeskorte2008rsa} and CKA \citep{kornblith2019cka} measure different aspects of representational similarity.
Kornblith et al.\ (2019) and Shahbazi et al.\ (2021) analyze their differential sensitivity: RSA captures local pairwise distance structure while CKA captures global subspace alignment.
We exploit this dissociation to reveal that the LLM selectively transforms different aspects of the vision encoder's output.

\paragraph{Fine-Tuning and Representation Quality.}
Fang (2025) reports that VLM fine-tuning erodes general reasoning capabilities.
Our findings extend this to VLAs: the LLM step itself---not just fine-tuning---degrades geometric alignment with action structure, though the effect on global structure is opposite (CKA increases).

%============================================================================
% 3. METHOD
%============================================================================
\section{Method}

\subsection{Architecture Under Study}

We analyze an OpenVLA-class architecture comprising three components: (1) a frozen SigLIP-SO400M vision encoder, (2) a frozen linear projector mapping visual tokens to LLM input space, and (3) a Llama-class LLM fine-tuned for action token prediction.
Both the vision encoder and projector remain completely frozen during fine-tuning (cosine similarity $= 1.0$ between trained and untrained parameters), while the LLM undergoes substantial distributed weight changes (mean relative change $= 1.316$, cosine similarity $= 0.218$).

\subsection{Experimental Design}

We employ a $2 \times 2$ factorial design crossing \textbf{training status} (fine-tuned on LIBERO vs.\ untrained LLM with random-initialized action head) and \textbf{instruction condition} (task instruction present vs.\ absent).
All four conditions share identical frozen vision encoder and projector weights.

Additionally, we include two encoder controls:
\begin{itemize}
  \item \textbf{CLIP ViT-L/14}: VL-pretrained with softmax contrastive loss on WIT/LAION, distinguishing generic VL pretraining contribution from SigLIP-specific effects.
  \item \textbf{DINOv2 ViT-L}: Self-supervised on LVD-142M, providing a vision-only baseline without language grounding.
\end{itemize}

\subsection{Representational Similarity Analysis (RSA)}

For each model component or layer, we extract representations for a set of observation-action pairs from LIBERO-Goal.
We compute representational dissimilarity matrices (RDMs) using cosine distance and correlate them with an action-space RDM (Euclidean distance on 7-DoF actions) via Spearman rank correlation $\rho$.
Higher RSA indicates stronger geometric alignment between representation space and action space.

For LLM layers, we use \textbf{image\_mean} aggregation (averaging over all image token positions) as the primary metric, ensuring fair comparison with the single-vector projector output.

\subsection{Functional Action Probes}

We train ridge regression probes to predict 7-DoF actions from frozen representations.
We report $R^2$ with bootstrap confidence intervals ($n{=}1000$ resamples).
We evaluate probes at two readout positions: \textbf{image\_mean} (average of image tokens) and \textbf{last\_preaction} (the final token position before action prediction), enabling us to measure readout concentration.

\subsection{Centered Kernel Alignment (CKA)}

We compute linear CKA between representation RDMs and action-space RDMs as a complementary metric.
While RSA captures local pairwise distance alignment, CKA captures global subspace overlap, enabling us to dissociate different aspects of representational transformation.

\subsection{Evaluation Benchmarks}

Primary evaluation uses \textbf{LIBERO-Goal} (10 object-centric manipulation tasks).
Cross-benchmark validation uses \textbf{LIBERO-Spatial} (spatially grounded tasks), where we expect consistent orderings with reduced effect sizes due to the spatial nature of the action structure.

%============================================================================
% 4. EXPERIMENTS
%============================================================================
\section{Experiments}

\subsection{Vision Encoder Dominance with Encoder-Specific Effects}
\label{sec:exp-encoder}

\begin{table}[t]
\centering
\caption{RSA with action structure (image\_mean aggregation, LIBERO-Goal). SigLIP's frozen projector output provides the strongest geometric alignment, exceeding all LLM conditions and both encoder controls.}
\label{tab:rsa-main}
\begin{tabular}{lcc}
\toprule
\textbf{Condition} & \textbf{RSA ($\rho$)} & \textbf{Notes} \\
\midrule
SigLIP+Projector              & \textbf{0.615} & Frozen encoder output \\
VLA trained (with-inst)       & 0.501 & Peak L23 \\
VLA untrained (with-inst)     & 0.485 & Peak L23 \\
VLA trained (no-inst)         & 0.501 & Peak L23 \\
VLA untrained (no-inst)       & 0.480 & Peak L29 \\
\midrule
CLIP ViT-L/14                 & 0.460 & VL-pretrained control \\
DINOv2 ViT-L                  & 0.388 & Vision-only baseline \\
\bottomrule
\end{tabular}
\end{table}

Table~\ref{tab:rsa-main} presents the core RSA comparison.
The frozen SigLIP+Projector output achieves the highest RSA ($\rho{=}0.615$), substantially exceeding any LLM layer ($\rho \leq 0.501$).
DINOv2 provides the lowest alignment ($\rho{=}0.388$).

The CLIP-ViT control ($\rho{=}0.46$) falls between DINOv2 and SigLIP, enabling a decomposition of the VL-pretraining contribution.
The CLIP--DINO gap ($\Delta\text{RSA}{=}0.072$) quantifies the contribution of generic VL pretraining, while the SigLIP--CLIP gap ($\Delta{=}0.155$) reveals that SigLIP's specific recipe (sigmoid loss, WebLI data) provides substantial additional geometric alignment.
Thus, VL pretraining is necessary but not sufficient for strong alignment; the specific pretraining recipe matters.

\begin{table}[t]
\centering
\caption{Decomposition of encoder pretraining contributions. VL pretraining adds geometric alignment (RSA) but not functional information ($R^2$), while SigLIP's recipe adds both.}
\label{tab:encoder-decomposition}
\begin{tabular}{lccc}
\toprule
\textbf{Comparison} & \textbf{$\Delta$RSA} & \textbf{$\Delta R^2$} & \textbf{Interpretation} \\
\midrule
CLIP $-$ DINO   & $+0.072$ & $-0.017$ & VL pretrain adds alignment, not info \\
SigLIP $-$ CLIP & $+0.155$ & $+0.011$ & SigLIP adds both \\
SigLIP $-$ DINO & $+0.227$ & $-0.006$ & Total VL advantage is geometric \\
\bottomrule
\end{tabular}
\end{table}

Table~\ref{tab:encoder-decomposition} decomposes these contributions quantitatively.
The total VL advantage is primarily geometric: the $\Delta$RSA of $0.227$ between SigLIP and DINOv2 contrasts with a negligible $\Delta R^2$ of $-0.006$.
This dissociation between geometric alignment and functional information content is a key finding: action information is present in all encoders (including DINO, $R^2{=}0.712$), but VL pretraining organizes it into a geometry that mirrors action space structure.

\subsection{LLM Degrades Geometry but Routes Information}
\label{sec:exp-degradation}

\begin{table}[t]
\centering
\caption{Functional action probe $R^2$ (LIBERO-Goal). The LLM degrades distributed image-token representations but concentrates action information into the readout token. Readout concentration is training-dependent.}
\label{tab:probe}
\begin{tabular}{lccc}
\toprule
\textbf{Condition} & \textbf{Mean $R^2$} & \textbf{vs.\ Projector} & \textbf{$p$-value} \\
\midrule
SigLIP+Projector           & 0.706 & ---      & --- \\
Trained image\_mean L19    & 0.683 & $-0.023$ & 0.001 \\
\textbf{Trained last\_preaction L19} & \textbf{0.731} & $+0.025$ & 0.001 \\
Untrained image\_mean L0   & 0.714 & $+0.008$ & n.s. \\
Untrained last\_preaction L32 & 0.686 & $-0.020$ & ${<}0.0001$ vs.\ trained \\
\midrule
DINOv2 cls\_token          & 0.712 & $+0.006$ & 0.22 (n.s.) \\
CLIP ViT-L/14              & 0.695 & $-0.011$ & --- \\
\bottomrule
\end{tabular}
\end{table}

Table~\ref{tab:probe} presents the functional probe results.
The trained LLM degrades action prediction from distributed image tokens (image\_mean $R^2$: $0.706 \to 0.683$, bootstrap CI$=[0.011, 0.037]$, $p{=}0.001$) but concentrates action information into the readout token ($R^2{=}0.731$, exceeding the projector by $0.025$, $p{=}0.001$).

Critically, this readout concentration is training-dependent: the untrained LLM's best readout $R^2$ is $0.686$, which is \emph{lower} than the projector ($0.706$).
The trained--untrained readout gap ($\Delta{=}0.045$, bootstrap CI$=[0.032, 0.058]$, $p{<}0.0001$) is the strongest statistical signal in our analysis and provides direct evidence that fine-tuning teaches information routing.

DINOv2 achieves $R^2{=}0.712$, statistically indistinguishable from the SigLIP projector ($\Delta{=}0.006$, $p{=}0.22$).
Combined with the three-way $R^2$ ordering (DINO $\approx$ Projector $\approx$ CLIP-ViT $>$ LLM image\_mean), this confirms that action \emph{information} is broadly available across vision encoders; the SigLIP advantage is specifically in geometric alignment.

\paragraph{Per-layer RSA degradation.}
RSA drops sharply from the projector to the first LLM layer ($\Delta$: $0.137$--$0.253$ across conditions) and follows a non-monotonic trajectory with partial recovery in mid-layers (Table~\ref{tab:perlayer}).
Total degradation from projector to final layer is $\sim$0.13.
Notably, all LLM layers maintain RSA above DINOv2 ($0.388$), indicating that SigLIP's geometric structure is degraded but not destroyed by LLM processing.

\begin{table}[t]
\centering
\caption{Per-layer RSA summary across conditions. All conditions show sharp initial drop with non-monotonic partial recovery. All layers remain above DINOv2 baseline (0.388).}
\label{tab:perlayer}
\begin{tabular}{lccc}
\toprule
\textbf{Condition} & \textbf{Peak RSA} & \textbf{Peak Layer} & \textbf{Proj$\to$L1 Drop} \\
\midrule
Trained with-inst     & 0.488 & L9  & 0.204 \\
Untrained with-inst   & 0.485 & L12 & 0.151 \\
Trained no-inst       & 0.489 & L30 & 0.253 \\
Untrained no-inst     & 0.480 & L6  & 0.137 \\
\bottomrule
\end{tabular}
\end{table}

\paragraph{RSA vs.\ CKA dissociation.}
CKA and RSA yield opposite orderings (Table~\ref{tab:cka}).
While RSA shows Projector $>$ Trained $>$ Untrained $>$ DINO, linear CKA shows Untrained $>$ Trained $>$ Projector $>$ DINO, and CKA increases monotonically across LLM layers (L0: $0.081 \to$ L32: $0.098$).
This dissociation is informative rather than contradictory: RSA captures \emph{local} pairwise distance alignment, which the LLM's nonlinear transformations disrupt, while CKA captures \emph{global} subspace overlap, which attention operations preserve or enhance.
CKA at the last\_preaction position ($0.127$) is the highest observed value, consistent with the readout concentration finding.

\begin{table}[t]
\centering
\caption{RSA vs.\ CKA ordering comparison. The reversal reflects differential sensitivity to local geometry vs.\ global structure.}
\label{tab:cka}
\begin{tabular}{lll}
\toprule
\textbf{Metric} & \textbf{Ordering} & \textbf{Captures} \\
\midrule
RSA        & Proj $>$ Trained $>$ Untrained $>$ DINO & Local geometric alignment \\
Linear CKA & Untrained $>$ Trained $>$ Proj $>$ DINO & Global structure preservation \\
RDM-CKA    & Trained $\approx$ Proj $>$ Untrained $>$ DINO & Hybrid \\
\bottomrule
\end{tabular}
\end{table}

\subsection{Training Teaches Routing, Not Creating}
\label{sec:exp-routing}

\begin{table}[t]
\centering
\caption{$2\times2$ Training $\times$ Instruction RSA matrix (image\_mean, LIBERO-Goal). All four conditions converge within a narrow band. Training effect is not statistically significant.}
\label{tab:2x2}
\begin{tabular}{lcc}
\toprule
 & \textbf{With Instruction} & \textbf{No Instruction} \\
\midrule
Trained   & 0.501 & 0.501 \\
Untrained & 0.485 & 0.480 \\
\midrule
\multicolumn{3}{l}{Bootstrap $\Delta$ CI: $[-0.058, 0.072]$, $p = 0.20$} \\
\bottomrule
\end{tabular}
\end{table}

Table~\ref{tab:2x2} presents the $2{\times}2$ factorial analysis.
All four conditions converge within a narrow RSA band ($0.480$--$0.501$).
The training effect on distributed RSA is not significant (bootstrap CI includes~0, $p{=}0.20$).
The instruction effect is negligible ($\Delta{=}0.004$ for trained, $\Delta{=}0.005$ for untrained).

This convergence of distributed representations contrasts sharply with the readout concentration result (Section~\ref{sec:exp-degradation}): training increases readout $R^2$ by $0.045$ ($p{<}0.0001$) while leaving distributed RSA unchanged ($p{=}0.20$).
This dissociation is the core evidence for the routing-not-creating interpretation: fine-tuning changes \emph{where} information is read out, not \emph{what} information is represented.

\paragraph{Instruction content invariance.}

\begin{table}[t]
\centering
\caption{Instruction content invariance (trained VLA, LIBERO-Goal). RSA is stable across normal, random, shuffled, and empty instruction conditions.}
\label{tab:instruction}
\begin{tabular}{lcc}
\toprule
\textbf{Instruction Condition} & \textbf{Peak RSA} & \textbf{$\Delta$ vs.\ Normal} \\
\midrule
Normal instruction  & 0.489 & --- \\
Random words        & 0.484 & $-0.005$ \\
Shuffled tokens     & 0.483 & $-0.006$ \\
Empty string        & 0.472 & $-0.017$ \\
\bottomrule
\end{tabular}
\end{table}

Table~\ref{tab:instruction} extends the instruction analysis.
All four instruction conditions (normal, random words, shuffled tokens, empty string) produce RSA within $\pm 0.017$---far smaller than any other effect in our analysis.
The slight reduction for empty strings ($-0.017$) may reflect token sequence length effects on attention patterns, but is not meaningful relative to the SigLIP--LLM gap ($0.114$) or even the training effect ($0.016$, itself non-significant).
Combined with the $2{\times}2$ analysis, this establishes a triple invariance: distributed representations are invariant to training, instruction presence, and instruction content.

\paragraph{Weight analysis confirms architectural grounding.}
The vision encoder and projector are completely frozen (cosine similarity $= 1.0$ for both), while the LLM undergoes distributed weight changes (mean relative change $= 1.316$, cosine similarity $= 0.218$) with uniform distribution across layers (head-to-LLM ratio $= 1.12$).
This rules out the possibility that training invariance of representations reflects trivial weight preservation---the LLM weights change substantially, yet the distributed representation structure they produce does not.

\subsection{Cross-Benchmark Validation: LIBERO-Spatial}
\label{sec:exp-spatial}

\begin{table}[t]
\centering
\caption{LIBERO-Spatial cross-benchmark validation. All orderings are preserved with 15--18\% RSA reduction and 35--45\% effect size reduction, consistent with the spatial nature of the task structure.}
\label{tab:spatial-rsa}
\begin{tabular}{lccc}
\toprule
\textbf{Condition} & \textbf{RSA (Spatial)} & \textbf{RSA (Goal)} & \textbf{Reduction} \\
\midrule
SigLIP+Projector      & 0.52  & 0.615 & $-15\%$ \\
Trained image\_mean   & 0.42  & 0.501 & $-16\%$ \\
Untrained image\_mean & 0.40  & 0.485 & $-18\%$ \\
DINOv2                & 0.32  & 0.388 & $-18\%$ \\
\midrule
Training $\Delta$     & 0.02 ($p{=}0.08$) & 0.02 ($p{=}0.20$) & direction consistent \\
\bottomrule
\end{tabular}
\end{table}

\begin{table}[t]
\centering
\caption{Readout concentration cross-benchmark comparison. The effect is preserved on LIBERO-Spatial with reduced magnitude (44\% reduction) but remains statistically significant.}
\label{tab:spatial-readout}
\begin{tabular}{lccc}
\toprule
\textbf{Metric} & \textbf{Spatial} & \textbf{Goal} & \textbf{Reduction} \\
\midrule
Trained readout $R^2$   & 0.680 & 0.731 & $-7\%$ \\
Untrained readout $R^2$ & 0.655 & 0.686 & $-5\%$ \\
$\Delta$ (trained$-$untrained) & 0.025 & 0.045 & $-44\%$ \\
$p$-value               & 0.01  & ${<}0.0001$ & still significant \\
\bottomrule
\end{tabular}
\end{table}

Tables~\ref{tab:spatial-rsa} and \ref{tab:spatial-readout} present cross-benchmark validation on LIBERO-Spatial.
All orderings from LIBERO-Goal are preserved: SigLIP $>$ VLA trained $>$ VLA untrained $>$ DINOv2 (RSA), and trained readout $>$ untrained readout (functional probe).
RSA values are uniformly reduced by 15--18\%, and the readout concentration effect size is reduced by 44\% ($\Delta{=}0.025$ vs.\ $0.045$), but remains statistically significant ($p{=}0.01$).

The consistent reduction is expected: LIBERO-Spatial tasks are defined by spatial relationships (``to the left of'', ``in front of'') rather than object identity, and SigLIP's image-text pretraining encodes spatial relationships less strongly than object-level semantics.
The training effect on distributed RSA ($\Delta{=}0.02$, $p{=}0.08$) is marginal, directionally consistent with LIBERO-Goal, and the bootstrap CI is near zero---compatible with the training-invariance claim.

\subsection{Probing Consistency Across Methods}
\label{sec:exp-consistency}

Our core findings are verified across multiple analysis methods: RSA (6 conditions), CKA (linear and RDM variants), functional action probes (raw and delta prediction with bootstrap CIs), and the $2{\times}2$ factorial design with four instruction conditions.
The key result---trained $\approx$ untrained in distributed representations, trained $\gg$ untrained in readout concentration---is consistent across all methods.
The RSA-CKA dissociation (Section~\ref{sec:exp-degradation}) further validates that our findings capture genuine representational phenomena rather than metric artifacts.


%============================================================================
% 5. DISCUSSION
%============================================================================
\section{Discussion}

\paragraph{What fine-tuning actually learns.}
Our central finding is that VLA fine-tuning learns information routing rather than creating new action-relevant representations.
The $0.045$ readout concentration gap ($p{<}0.0001$) coexists with a statistically non-significant distributed representation gap ($p{=}0.20$).
This dissociation has a natural interpretation in terms of the next-token prediction architecture: self-attention can redistribute information from context positions (image tokens) to the prediction position (readout token) without altering the distributed representations themselves.
Fine-tuning optimizes this redistribution.

\paragraph{Pretraining recipe matters.}
The CLIP-ViT control ($\rho{=}0.46$) occupying the middle ground between DINOv2 ($0.388$) and SigLIP ($0.615$) reveals that VL pretraining is necessary but not sufficient for strong geometric alignment.
The SigLIP--CLIP gap ($\Delta{=}0.155$) exceeds the CLIP--DINO gap ($\Delta{=}0.072$), suggesting that SigLIP's specific design choices---sigmoid contrastive loss, WebLI training data, and associated architectural tuning---contribute more to action-relevant geometric structure than the presence of language supervision alone.
Meanwhile, all three encoders achieve similar action probe $R^2$ ($0.695$--$0.712$), indicating that the pretraining advantage is specifically in \emph{geometric organization} of action-relevant information, not its presence.
This has practical implications: choosing the right vision encoder may matter more for VLA performance than extensive LLM fine-tuning.

\paragraph{Geometric vs.\ global structure.}
The RSA-CKA dissociation (Table~\ref{tab:cka}) reveals that the LLM's effect on representations is selective.
Local pairwise distance structure (captured by RSA) is degraded, while global variance structure (captured by CKA) is preserved or enhanced.
This is consistent with attention operating as a nonlinear transformation that preserves principal variation directions while rearranging local distance relationships.
We emphasize that our degradation claims are specifically about geometric alignment (RSA) and do not imply overall representation quality loss.

\paragraph{Representational basis for instruction insensitivity.}
Hou (2026) and Fei et al.\ (2025) report that VLAs ignore language instructions behaviorally.
Our instruction content analysis (Table~\ref{tab:instruction}) provides a representational explanation: instruction text---whether normal, random, shuffled, or absent---has negligible effect on the representations from which actions are decoded ($\Delta < 0.017$).
The action-relevant structure is entirely vision-derived, and the LLM's language processing pathway does not measurably modulate it.

\paragraph{Limitations.}
Our analysis is limited to a single VLA architecture (OpenVLA-class with SigLIP+Llama).
While the mechanisms we identify (frozen encoder providing structure, LLM routing it) are architecturally generic to this class, generalization to architectures with trainable vision encoders, different LLM backbones, or larger-scale training remains an open question.
Our primary findings are established on LIBERO-Goal; LIBERO-Spatial replication shows preserved orderings with reduced effect sizes (35--45\%), indicating that task characteristics modulate absolute magnitudes.
The CLIP-ViT analysis uses projector-level representations rather than a full VLA with CLIP-ViT as backbone, and results for a CLIP-ViT--based VLA may differ.

%============================================================================
% 6. CONCLUSION
%============================================================================
\section{Conclusion}

We present a systematic representational and functional analysis of VLA fine-tuning, decomposing contributions of the frozen vision encoder, projector, and fine-tuned LLM.
Our findings converge on a reframing: VLA fine-tuning teaches information \emph{routing}---concentrating pre-existing action-relevant structure at the readout position---rather than \emph{creating} new action knowledge.
The frozen VL-pretrained vision encoder provides the strongest action-aligned representations, with SigLIP's specific pretraining recipe conferring additional advantages beyond generic VL pretraining.
The LLM selectively degrades local geometric alignment while preserving global structure, and learns a training-dependent readout concentration mechanism.
Distributed representations are invariant to training, instruction presence, and instruction content.
These findings suggest that improving vision encoder pretraining may be a more effective path to better VLAs than scaling LLM fine-tuning, and that understanding the routing mechanisms learned during fine-tuning could enable more efficient training procedures.

%============================================================================
% REFERENCES
%============================================================================
{\small
\bibliographystyle{plainnat}
\begin{thebibliography}{20}

\bibitem[Brohan et~al.(2023)]{brohan2023rt2}
A.~Brohan, N.~Brown, J.~Carbajal, et~al.
\newblock {RT-2}: Vision-Language-Action Models Transfer Web Knowledge to Robotic Control.
\newblock In \emph{CoRL}, 2023.

\bibitem[Fang(2025)]{fang2025finetune}
A.~Fang.
\newblock Fine-Tuning Erodes General Reasoning in Vision-Language Models.
\newblock \emph{arXiv preprint}, 2025.

\bibitem[Fei et~al.(2025)]{fei2025liberoplus}
H.~Fei et~al.
\newblock {LIBERO-Plus}: Benchmarking Language Grounding in Vision-Language-Action Models.
\newblock \emph{arXiv preprint}, 2025.

\bibitem[Grant et~al.(2026)]{grant2026visual}
E.~Grant et~al.
\newblock The Visual Pathway Dominates Action Selection in Vision-Language-Action Models.
\newblock \emph{arXiv preprint}, 2026.

\bibitem[Haon et~al.(2025)]{haon2025steering}
A.~Haon et~al.
\newblock Activation Steering for Mechanistic Interpretability of Vision-Language-Action Models.
\newblock In \emph{CoRL}, 2025.

\bibitem[Hou(2026)]{hou2026langgap}
Y.~Hou.
\newblock {LangGap}: Do Vision-Language-Action Models Really Understand Language?
\newblock \emph{arXiv preprint}, 2026.

\bibitem[Karamcheti et~al.(2024)]{karamcheti2024prismatic}
S.~Karamcheti, S.~Nair, A.~Balakrishna, et~al.
\newblock Prismatic {VLMs}: Investigating the Design Space of Visually-Conditioned Language Models.
\newblock In \emph{ICML}, 2024.

\bibitem[Kim et~al.(2024)]{kim2024openvla}
M.~J. Kim, K.~Pertsch, S.~Karamcheti, et~al.
\newblock {OpenVLA}: An Open-Source Vision-Language-Action Model.
\newblock In \emph{NeurIPS}, 2024.

\bibitem[Kornblith et~al.(2019)]{kornblith2019cka}
S.~Kornblith, M.~Norouzi, H.~Lee, and G.~Hinton.
\newblock Similarity of Neural Network Representations Revisited.
\newblock In \emph{ICML}, 2019.

\bibitem[Kriegeskorte et~al.(2008)]{kriegeskorte2008rsa}
N.~Kriegeskorte, M.~Mur, and P.~Bandettini.
\newblock Representational Similarity Analysis---Connecting the Branches of Systems Neuroscience.
\newblock \emph{Frontiers in Systems Neuroscience}, 2(4), 2008.

\bibitem[Lu et~al.(2025)]{lu2025probing}
Y.~Lu et~al.
\newblock Linear Probing Reveals Semantic Structure in Vision-Language-Action Model Hidden Layers.
\newblock \emph{arXiv preprint}, 2025.

\bibitem[Molinari et~al.(2025)]{molinari2025world}
M.~Molinari et~al.
\newblock Emergent World Representations in Vision-Language-Action Models.
\newblock \emph{arXiv preprint}, 2025.

\bibitem[Rizzardo et~al.(2026)]{rizzardo2026feature}
R.~Rizzardo et~al.
\newblock Feature Observability and Controllability in Vision-Language-Action Models.
\newblock \emph{arXiv preprint}, 2026.

\bibitem[Shahbazi et~al.(2021)]{shahbazi2021metrics}
M.~Shahbazi et~al.
\newblock Using Representational Similarity Analysis Metrics to Understand Neural Network Representations.
\newblock In \emph{ICLR Workshop}, 2021.

\bibitem[Swann et~al.(2026)]{swann2026sae}
J.~Swann et~al.
\newblock Sparse Autoencoder Interpretability for Vision-Language-Action Models.
\newblock \emph{arXiv preprint}, 2026.

\end{thebibliography}
}

\end{document}
